\documentclass[11pt]{article}

\usepackage[a4paper,margin=1in]{geometry}
\usepackage{amsmath,amssymb,amsthm,mathtools}
\usepackage{graphicx}
\usepackage{float}
\usepackage{microtype}
\usepackage[hidelinks]{hyperref}

\newtheorem{theorem}{Theorem}
\newtheorem{lemma}[theorem]{Lemma}
\newtheorem{remark}[theorem]{Remark}

\newcommand{\R}{\mathbb R}
\newcommand{\Z}{\mathbb Z}
\newcommand{\cG}{\mathcal G}
\newcommand{\cZ}{\mathcal Z}
\newcommand{\T}{T}
\newcommand{\M}{M}
\newcommand{\Dminus}{D^-}
\newcommand{\Dplus}{D^+}

\title{An odd-window Gabor frame on a non-lattice set of density two}
\author{Candidate counterexample packet}
\date{9 August 2026}

\begin{document}
\maketitle

\begin{center}
\textbf{Status: candidate counterexample, likely valid pending expert review.}
\end{center}

\begin{abstract}
Faulhuber asks whether the Lyubarskii--Nes non-frame theorem for odd
windows and lattices of density $(n+1)/n$ extends to arbitrary relatively
separated discrete sets of the same lower Beurling density.  We answer the
proposed universal extension negatively already for $n=1$.  An explicit odd
$C_c^\infty$ window $g$ and the non-lattice periodic set
\[
 \Gamma=\Z^2\cup\bigl((1/4,0)+\Z^2\bigr)
\]
give a Gabor frame for $L^2(\R)$, while
$\Dminus(\Gamma)=\Dplus(\Gamma)=2$.  The proof is exact: the Zak transform
of $g$ vanishes on two vertical lines, and the quarter-translate moves those
lines to a disjoint pair.  The two-window Zak multiplier is consequently
strictly positive on the compact torus.
\end{abstract}

\section{Source question}

The source is M.~Faulhuber, \emph{On the Parity under Metaplectic Operators
and an Extension of a Result of Lyubarskii and Nes}, arXiv:1901.01220v2
\cite{Faulhuber2019}.  On PDF page 9, after extending the odd-window
non-frame theorem from separable lattices to all lattices in dimension one,
the paper asks:

\begin{quote}
Does the result in Corollary 1.2 depend on the group structure of $\Lambda$
or does the result hold for any relatively separated, discrete set
$\Gamma\subset\R^2$ with lower Beurling density $(n+1)/n$?
\end{quote}

\begin{figure}[H]
  \centering
  \includegraphics[width=0.97\textwidth]{figures/open_problem_crop.png}
  \caption{The open questions in arXiv:1901.01220v2, PDF page 9.}
\end{figure}

The question is naturally universal in $n$ and in the odd window.  The
construction below refutes that universal extension at $n=1$.  It does not
settle analogous assertions for each fixed $n\geq2$, nor either of the two
higher-dimensional questions in the source.

\section{Definitions and statement}

For $x,\omega\in\R$, set
\[
  \T_xf(t)=f(t-x),\qquad
  \M_\omega f(t)=e^{2\pi i\omega t}f(t).
\]
For a discrete $\Gamma\subset\R^2$, the associated Gabor system is
\[
 \cG(g,\Gamma)=\{\M_\omega\T_xg:(x,\omega)\in\Gamma\}.
\]
It is a frame if there are $0<A\leq B<\infty$ such that
\[
 A\|f\|_2^2\leq
 \sum_{(x,\omega)\in\Gamma}|\langle f,\M_\omega\T_xg\rangle|^2
 \leq B\|f\|_2^2
 \quad(f\in L^2(\R)).
\]

Define
\begin{equation}\label{eq:window}
 g(t)=
 \begin{cases}
 t\exp\!\left(-\dfrac{1}{1-4t^2}\right),& |t|<\dfrac12,\\[6pt]
 0,& |t|\geq\dfrac12,
 \end{cases}
\end{equation}
and
\begin{equation}\label{eq:set}
 \Gamma=\Z^2\cup\bigl((1/4,0)+\Z^2\bigr).
\end{equation}

\begin{theorem}[Counterexample to the irregular-set extension]
\label{thm:main}
The function $g$ in \eqref{eq:window} is odd and belongs to Feichtinger's
algebra $S_0(\R)$.  The set $\Gamma$ in \eqref{eq:set} is relatively
separated and discrete, is not a lattice, and satisfies
\[
 \Dminus(\Gamma)=\Dplus(\Gamma)=2.
\]
Nevertheless, $\cG(g,\Gamma)$ is a frame for $L^2(\R)$.  Hence the first
open question above has a negative answer as a universal extension: the
Lyubarskii--Nes obstruction does depend on more than lower Beurling density.
\end{theorem}

\section{Proof intuition}

At the critical lattice $\Z^2$, the Zak transform turns a multi-window
Gabor frame operator into multiplication by the sum of the squared Zak
transforms of its windows.  Thus two windows give a frame exactly when their
Zak transforms never vanish simultaneously.

The compact support of \eqref{eq:window} makes its Zak transform completely
transparent.  On one half of a fundamental interval only the translate
$k=0$ contributes, and on the other half only $k=1$ contributes.  Since $g$
vanishes only at its center and endpoints, its Zak transform has precisely
the vertical zero lines $x=0$ and $x=1/2$ modulo one.  Translating $g$ by
$1/4$ moves these to $x=1/4$ and $x=3/4$.  The two pairs are disjoint, so
the sum of squares has a positive minimum.  The two windows over $\Z^2$
are exactly the single-window Gabor system over the two-coset set $\Gamma$.

\section{Zak-transform lemma}

Use the convention
\begin{equation}\label{eq:zak}
 \cZ f(x,\omega)=\sum_{k\in\Z}f(x-k)e^{2\pi i k\omega}.
\end{equation}
The Zak transform extends to a unitary map from $L^2(\R)$ onto the usual
quasi-periodic realization of $L^2([0,1]^2)$.

\begin{lemma}[Critical multi-window criterion]\label{lem:zak}
Let $h_1,\ldots,h_m\in S_0(\R)$.  Then
\[
 \bigcup_{j=1}^m\cG(h_j,\Z^2)
\]
is a frame if and only if there are $0<A\leq B<\infty$ such that
\begin{equation}\label{eq:multiplier}
 A\leq\sum_{j=1}^m|\cZ h_j(x,\omega)|^2\leq B
\end{equation}
on $[0,1]^2$.  In particular, because these Zak transforms are continuous,
the union is a frame if they have no common zero.
\end{lemma}

\begin{proof}
Under $\cZ$, the functions $\M_\ell\T_kh_j$ become, up to a unimodular
factor,
\[
 e^{2\pi i(k\omega-\ell x)}\cZ h_j(x,\omega).
\]
Parseval's identity for the Fourier basis
$\{e^{2\pi i(k\omega-\ell x)}\}_{k,\ell\in\Z}$ shows that the frame
operator is unitarily equivalent to multiplication by the function in
\eqref{eq:multiplier}.  Such a multiplication operator is bounded and
invertible exactly when its multiplier is bounded above and away from zero.
The continuity assertion follows from $S_0(\R)\subset W_0(\R)$; alternatively
it is immediate for the particular compactly supported smooth windows used
below.  This standard Zak diagonalization is also written out for periodic
configurations in \cite{Faulhuber2025}.
\end{proof}

\section{Proof of Theorem \ref{thm:main}}

\begin{proof}
The bump in \eqref{eq:window} is smooth at $\pm1/2$ with all derivatives
zero there.  Hence $g\in C_c^\infty(\R)\subset\mathcal S(\R)\subset
S_0(\R)$.  It is odd, and
\begin{equation}\label{eq:nonzero}
 g(t)\neq0\quad\text{for }0<|t|<1/2.
\end{equation}

We first determine the exact zero set of $\cZ g$ in a fundamental domain.
If $0<x<1/2$, only $k=0$ contributes to \eqref{eq:zak}, whereas if
$1/2<x<1$, only $k=1$ contributes.  Therefore
\begin{equation}\label{eq:piecewise}
 \cZ g(x,\omega)=
 \begin{cases}
 g(x),&0<x<1/2,\\
 e^{2\pi i\omega}g(x-1),&1/2<x<1.
 \end{cases}
\end{equation}
At $x=0$ the only possible interior contribution is $g(0)=0$, and at
$x=1/2$ the two endpoint contributions vanish.  Combining this with
\eqref{eq:nonzero} gives
\begin{equation}\label{eq:zeros}
 \{(x,\omega)\in[0,1)^2:\cZ g(x,\omega)=0\}
 =\bigl(\{0,1/2\}\times[0,1)\bigr).
\end{equation}

Translation by $a$ acts on the Zak transform by
\[
 \cZ(\T_a g)(x,\omega)=\cZ g(x-a,\omega).
\]
For $a=1/4$, the zero set modulo one is consequently
\[
 \{1/4,3/4\}\times[0,1),
\]
which is disjoint from \eqref{eq:zeros}.  Thus the continuous function
\[
 Q(x,\omega)=|\cZ g(x,\omega)|^2+
 |\cZ(\T_{1/4}g)(x,\omega)|^2
\]
is strictly positive on the compact torus.  It has a positive minimum and a
finite maximum.  Lemma \ref{lem:zak} proves that
\[
 \cG(g,\Z^2)\cup\cG(\T_{1/4}g,\Z^2)
\]
is a frame.  Since translations commute,
\[
 \M_\ell\T_k(\T_{1/4}g)=\M_\ell\T_{k+1/4}g,
\]
so this multi-window system is exactly $\cG(g,\Gamma)$.

It remains to check the geometry of $\Gamma$.  It is the union of two
distinct cosets of $\Z^2$, hence is uniformly discrete (the minimum
separation is $1/4$) and therefore relatively separated.  Every unit-area
fundamental cell contains two points up to boundary conventions, so standard
periodic counting gives
$\Dminus(\Gamma)=\Dplus(\Gamma)=2$.  Finally, $0,(1/4,0)\in\Gamma$ but
$(1/2,0)\notin\Gamma$; thus $\Gamma$ is not an additive subgroup and hence
is not a lattice.  All hypotheses in the source question at $n=1$ are met,
while the proposed non-frame conclusion fails.
\end{proof}

\section{Verification, scope, and novelty status}

The proof is analytic and does not depend on computation.  The accompanying
script samples the exact piecewise Zak multiplier, checks the four vertical
zero classes, and reports a positive numerical lower bound as a regression
test.  This is a sanity check, not part of the proof.

The counterexample settles the first source question negatively in its
universal form, already at $n=1$.  It does not decide whether a corresponding
non-frame theorem might hold at each fixed density $(n+1)/n$ for $n\geq2$,
and it does not address the higher-dimensional questions.

A bounded novelty search covered the run's ledger/indexes and arXiv searches
for the source title, ``odd Gabor window'', ``relatively separated'',
``lower Beurling density'', ``Lyubarskii--Nes'', compactly supported odd
windows, and unions of shifted integer lattices.  It also inspected the
closely related later paper arXiv:2502.09510 \cite{Faulhuber2025}.  That paper
uses the same standard Zak multiplier for periodic configurations and proves
generic frame statements under a finite separated-zero hypothesis, but the
search did not find the explicit window \eqref{eq:window}, the two-coset
counterexample \eqref{eq:set}, or an explicit answer to the 2019 question.
Novelty confidence is therefore moderate, not bibliographically certified;
the mechanism is elementary and close to known Zak-transform technology.

\paragraph{Human-review recommendation.}
Prioritize the identification of the source quantifiers, the critical
multi-window Zak criterion, and the exact endpoint zero calculation in
\eqref{eq:piecewise}--\eqref{eq:zeros}.  These are the only substantive
points on which the counterexample depends.

\begin{thebibliography}{9}

\bibitem{Faulhuber2019}
M.~Faulhuber,
\emph{On the Parity under Metaplectic Operators and an Extension of a
Result of Lyubarskii and Nes}, Results Math. 75 (2020), article 8;
arXiv:1901.01220v2.

\bibitem{LyubarskiiNes}
Y.~Lyubarskii and P.~G.~Nes,
\emph{Gabor frames with rational density}, J. Fourier Anal. Appl. 19
(2013), 805--822; arXiv:1108.2684.

\bibitem{Faulhuber2025}
M.~Faulhuber,
\emph{Gabor systems with Hermite functions of order $n$ and oversampling
greater than $n+1$ which are not frames}, Sampling Theory, Signal
Processing, and Data Analysis (2025); arXiv:2502.09510v2.

\end{thebibliography}

\end{document}
